%--------------------------
\begin{frame}{Question 1}

\begin{block}{}
Le déposant a le droit de désigner comme mandataire pendant la phase
internationale :

\begin{itemize}
  \item[$\circ$] a. une personne qui a le droit d’exercer auprès de n’importe quel office récepteur.
  \item[$\circ$] b. toute personne qui a le droit d’exercer auprès du Bureau international en sa qualité
  d’office récepteur.
  \item[$\bullet$] c. toute personne qui a le droit d’exercer auprès de l’office récepteur auprès duquel la
  demande internationale est déposée.
  \item[$\circ$] d. Aucune des réponses a, b ou c n’est correcte.
\end{itemize}

\end{block}

\end{frame}

\begin{frame}{Question 1}

\textbf{Base juridique :} \pctart{49} + \pctrule{90.1}.  

\begin{itemize}
  \item[$\circ$] \textbf{a.} \; Faux : le déposant ne peut pas désigner un mandataire habilité
  auprès de \emph{n’importe quel} office récepteur, mais uniquement auprès de l’office national
  où la demande est déposée (\pctrule{90}.1(a)).

  \item[$\circ$] \textbf{b.} \; Faux : un mandataire ne peut exercer auprès du Bureau international
  que si la demande est déposée auprès du BI agissant comme office récepteur (\pctrule{90}.1(a)).

  \item[$\bullet$] \textbf{c.} \; Correct : le déposant peut désigner une personne ayant le droit
  d’exercer auprès de l’office récepteur auprès duquel la demande internationale est déposée
  (\pctrule{90}.1(a)).

  \item[$\circ$] \textbf{d.} \; Faux puisque la réponse \textbf{c} est correcte.
\end{itemize}

\end{frame}